\documentclass[a4paper]{article}

\usepackage[pages=all, color=black, position={current page.south}, placement=bottom, scale=1, opacity=1, vshift=5mm]{background}
\SetBgContents{

}      % copyright

\usepackage[margin=1in]{geometry} % full-width

% AMS Packages
\usepackage{amsmath}
\usepackage{amsthm}
\usepackage{amssymb}

% Unicode
\usepackage[utf8]{inputenc}
\usepackage{hyperref}
\hypersetup{
	unicode,
%	colorlinks,
%	breaklinks,
%	urlcolor=cyan, 
%	linkcolor=blue, 
	pdfauthor={Author One, Author Two, Author Three},
	pdftitle={A simple article template},
	pdfsubject={A simple article template},
	pdfkeywords={article, template, simple},
	pdfproducer={LaTeX},
	pdfcreator={pdflatex}
}

% Vietnamese
%\usepackage{vntex}

% Natbib
\usepackage[sort&compress,numbers,square]{natbib}
\bibliographystyle{mplainnat}

% Theorem, Lemma, etc
\theoremstyle{plain}
\newtheorem{theorem}{Theorem}
\newtheorem{corollary}[theorem]{Corollary}
\newtheorem{lemma}[theorem]{Lemma}
\newtheorem{claim}{Claim}[theorem]
\newtheorem{axiom}[theorem]{Axiom}
\newtheorem{conjecture}[theorem]{Conjecture}
\newtheorem{fact}[theorem]{Fact}
\newtheorem{hypothesis}[theorem]{Hypothesis}
\newtheorem{assumption}[theorem]{Assumption}
\newtheorem{proposition}[theorem]{Proposition}
\newtheorem{criterion}[theorem]{Criterion}
\theoremstyle{definition}
\newtheorem{definition}[theorem]{Definition}
\newtheorem{example}[theorem]{Example}
\newtheorem{remark}[theorem]{Remark}
\newtheorem{problem}[theorem]{Problem}
\newtheorem{principle}[theorem]{Principle}

\usepackage{graphicx, color}
\graphicspath{{fig/}}

%\usepackage[linesnumbered,ruled,vlined,commentsnumbered]{algorithm2e} % use algorithm2e for typesetting algorithms
\usepackage{algorithm, algpseudocode} % use algorithm and algorithmicx for typesetting algorithms
\usepackage{mathrsfs} % for \mathscr command

\usepackage{lipsum}

% Author info
\pagebreak
\title{Translation}
\author{María Jesús Dueñas Recuero}

\date{

 \\ \texttt{MJesus.duenas1@alu.uclm.es}\\[2ex]%
%	\today
}

\begin{document}
	\maketitle
	\newpage
	\begin{abstract}
		Given the increase in globalisation, every day we are in contact with a greater number of cultures, which forces us, in a healthy way, to learn new languages.

	However, thanks to the development of technology and artificial intelligence, it is not necessarily necessary to study a language for years in order to understand it, given that translator websites and the translators themselves allow us to interact with it fluently.

	It is an opportunity to broaden our knowledge and understanding of other languages. As well as guaranteeing the correct use of the language for work and business, as they not only translate a language but also interpret emotions and feelings.


	\end{abstract}
	\newpage
	\tableofcontents
	\newpage
	\section{Introduction}
	\label{sec:intro}
	
	\lipsum[2]
	
	\subsection{Preliminaries}
	\label{sec:pre}
	
	\lipsum[3]
	
	\subsection{Previous Results}
	\label{sec:prev-results}
	
	Null graphs are discussed in \cite{HararyR74}
	The concept of ``internally stable set'' was used in \cite{Berge57, Berge58}.
	
	\begin{theorem}
		\label{thrm:1}
		\lipsum[4]
	\end{theorem}
	\begin{proof}
		content...
	\end{proof}

	\begin{corollary}
	\label{cor:1}
	
	\lipsum[5]
	\end{corollary}

	Unordered List (taken from Overleaf)
	\begin{itemize}
		\item The individual entries are indicated with a black dot, a so-called bullet.
		\item The text in the entries may be of any length.
	\end{itemize}

	Ordered List (taken from Overleaf)
	\begin{enumerate}
		\item The labels consists of sequential numbers.
		\item The numbers starts at 1 with every call to the enumerate environment.
	\end{enumerate}

	\begin{table}[ht]
		\centering
		\begin{tabular}{|c|c|}
			\hline
			\textbf{Odd} & \textbf{Even} \\
			\hline\hline
			One & Two \\
			\hline
			Three & Four \\
			\hline
		\end{tabular}
		\caption{This is a table}
		\label{tbl:1}
	\end{table}

	Table~\ref*{tbl:1} is an example of a table.
	
	\section{More Examples}
	\label{sec:examples}
	
	Now we include a figure.
	(See Figure~\ref{fig:example}.)
	\begin{figure}[ht]
		\centering
		\includegraphics[width=0.3\textwidth]{example}
		\caption{An example of a figure}
		\label{fig:example}
	\end{figure}
	
	\paragraph{Acknowledgements} \lipsum[6]
	
%	\newpage
	\bibliography{refs}
	
	\appendix
	
	\section{Omitted Proof in Section~\ref{sec:examples}}
	\label{app:1}
	
	\lipsum[7]
	
\end{document}